
\begin{center}
  \large
  \textbf{QUASINORMAL MODES IN SCHWARZSCHILD ANTI DE-SITTER SPACETIME}
\end{center}
\addcontentsline{toc}{chapter}{ABSTRACT}
\thispagestyle{empty}

\begin{flushleft}
  \setlength{\tabcolsep}{0pt}
  \bfseries
  \begin{tabular}{ll@{\hspace{6pt}}l}
  Student Name &:& Ilham Rasyid Machfudi\\
  NRP &:& 5001221124\\
  Department&:& Physics FSAD-ITS\\
  Advisor&:& Dr. rer. nat. Bintoro Anang Subagyo\\
  \end{tabular}
  \vspace{4ex}
\end{flushleft}
\textbf{Abstract}

% Isi Abstract
Quasinormal modes are the characteristic damped oscillations excited when a black
hole is subjected to small perturbations, governed by complex frequencies
$\omega = \omega_R - i\omega_I$ whose real part represents the oscillation
frequency and whose imaginary part represents the damping rate. This study
determines and compares the quasinormal mode frequency spectra of scalar field
perturbations in two spacetimes: the asymptotically flat Schwarzschild spacetime
and the Schwarzschild Anti-de Sitter spacetime. The eigenfrequencies are obtained
using the Frobenius method combined with the continued fraction method; for the
Schwarzschild Anti-de Sitter case, the resulting multi-term recurrence relation
is first reduced to a three-term relation through Gaussian elimination. The
computations are carried out for massless ($\mu = 0$) and massive ($\mu = 1$)
scalar fields, angular momenta $\ell = 0, 1, 2$, and both the fundamental mode
and the first overtone. The results show that in the Schwarzschild spacetime an
increase in the black hole mass lowers both the real and imaginary parts of the
frequency, whereas in the Schwarzschild Anti-de Sitter spacetime an increase in
the event horizon radius $r_+$ instead raises both components. An increase in
angular momentum consistently raises the real part while lowering the imaginary
part of the frequency in both spacetimes, whereas an increase in the overtone
index and the addition of scalar field mass produce differing trends between the
Schwarzschild and Schwarzschild Anti-de Sitter spacetimes. All obtained modes
have $\omega_I > 0$, indicating that both the Schwarzschild and Schwarzschild
Anti-de Sitter black holes are stable against scalar field perturbations, for
both massless and massive fields.
\vspace{2ex}\\
\noindent \textbf{Keywords: \emph{Quasinormal Modes, SAdS, Scalar Field, Schwarzschild}}